% answers_short.tex —— 2010 年·普通高等学校招生考试·江苏卷 · 简短答案速查
% 由 tools/make_short_answers.py 自 answers.tex 的 \gksolhead/\gkanswer 提取，
% 并与 2010/2010_数学真题_江苏_答案要点.md 逐题比对后生成，勿手改（改源 MD 后重跑）。

\begin{gknotice}
\textbf{使用说明}：本节只给最终答案便于快速核对，完整推导见解析版。答案为编写者独立推导，\textbf{非官方标准答案}。
\end{gknotice}

\gkpart{填空题}
\begin{itemize}[leftmargin=2.6em]
  \item[\textbf{1.}] $1$。
  \item[\textbf{2.}] $2$。
  \item[\textbf{3.}] $\dfrac12$。
  \item[\textbf{4.}] $30$。
  \item[\textbf{5.}] $-1$。
  \item[\textbf{6.}] $4$。
  \item[\textbf{7.}] $63$。
  \item[\textbf{8.}] $21$。
  \item[\textbf{9.}] $(-13,13)$。
  \item[\textbf{10.}] $\dfrac23$。
  \item[\textbf{11.}] $(-1,\sqrt2-1)$。
  \item[\textbf{12.}] $27$。
  \item[\textbf{13.}] $4$。
  \item[\textbf{14.}] $\dfrac{32\sqrt3}{3}$。
\end{itemize}

\gkpart{解答题}
\begin{itemize}[leftmargin=2.6em]
  \item[\textbf{15.}] （1）$2\sqrt{10}$，$4\sqrt2$；（2）$t=-\dfrac{11}5$。
  \item[\textbf{16.}] （1）证明；（2）$\sqrt2$。
  \item[\textbf{17.}] （1）$H=124$；（2）$d=55\sqrt5$。
  \item[\textbf{18.}] （1）$x=\dfrac92$；（2）$T\left(7,\dfrac{10}3\right)$；（3）定点 $(1,0)$。
  \item[\textbf{19.}] （1）$a_n=(2n-1)d^2$；（2）证明，$c_{\max}=\dfrac92$。
  \item[\textbf{20.}] （1）① 证明；② $b\leqslant2$ 时在 $(1,+\infty)$ 递增；$b>2$ 时在 $\left(1,\dfrac{b+\sqrt{b^2-4}}2\right)$ 递减，在 $\left(\dfrac{b+\sqrt{b^2-4}}2,+\infty\right)$ 递增；（2）$0<m<1$。
\end{itemize}

\gkpart{附加题}
\begin{itemize}[leftmargin=2.6em]
  \item[\textbf{21.}] A 证明；B $k=\pm2$；C $a=2$ 或 $-8$；D 证明。
  \item[\textbf{22.}] （1）分布列见下；（2）$\dfrac{512}{625}$。
  \item[\textbf{23.}] （1）证明；（2）证明。
\end{itemize}
